\documentclass[10pt,a4paper]{amsart}
\usepackage[margin=19mm]{geometry}
\usepackage{amsmath,amssymb,mathtools}
\usepackage[scr=rsfs,cal=euler]{mathalfa}
\usepackage{xfrac}
\usepackage[unicode,hidelinks,bookmarks=false]{hyperref}
\newcommand{\ie}{i.e.\ }
\newcommand{\cf}{cf.\ }
\newcommand{\resp}{respectively\ }
\newcommand{\Subsection}{\subsection}
\providecommand{\nobreakdash}{\nobreak\hbox{-}}
\setlength{\emergencystretch}{2em}
\multlinegap0pt
\usepackage{bm}
\usepackage{bbm}
\usepackage{dsfont}
\usepackage{xspace}
\usepackage{cancel}
\usepackage{mathtools}
\usepackage{amsthm}
\makeatletter
\@ifundefined{problem}{\newtheorem{problem}{Problem}}{}
\@ifundefined{proposition}{\newtheorem{proposition}{Proposition}}{}
\makeatother
\makeatletter
\newcommand{\Expect}{\mathbb{E}}
\def\R{\mathbb{R}}
\expandafter\ifx\csname romannum\endcsname\relax
\newenvironment{romannum}{\begin{list}{{\upshape (\roman{rmnum})}}{\usecounter{rmnum}
\setlength{\leftmargin}{14pt}
\setlength{\rightmargin}{8pt}
\setlength{\itemsep}{2pt}
\setlength{\itemindent}{-1pt}
}}{\end{list}}
\fi
\newcommand{\trace}{\text{Tr}}
\newcommand{\ud}{\,\mathrm{d}}
\makeatother
\title[A Time-Reversal Control Synthesis for Steering the State of ]{A Time-Reversal Control Synthesis for Steering the State of Stochastic Systems}
\author{Yuhang Mei, Amirhossein Taghvaei, Ali Pakniyat}
\date{Source: arXiv:2504.00238v2 (CC BY 4.0)}
\begin{document}
\maketitle

\noindent\emph{Source and licence.} A Time-Reversal Control Synthesis for Steering the State of Stochastic Systems. Version arXiv:2504.00238v2 (CC BY 4.0), distributed under the Creative Commons Attribution 4.0 International licence, as recorded in the arXiv licence metadata for this exact version and shown on the abstract page https://arxiv.org/abs/2504.00238v2. Full rights notice: licenses/arxiv-2504.00238-CC-BY-4.0.txt.\par\smallskip

\noindent\textit{Editorial scope.} Counted unit: the Proposition whose source label is \texttt{prop:terminal-error} in the article (source0.tex lines 1151--1158, source label \texttt{prop:terminal-error}), delivered together with the article's own complete original proof of it (source0.tex lines 1159--1177, one \texttt{proof} environment of 19 lines). No mathematics is written, repaired or re-derived: statement and proof are the article's own text line for line. Required context retained uncounted: prob:steer-to-y (lines 677--679); prob:approx-steer-to-y (lines 683--685); eq:kappa-identity (lines 788--790); eq:k-star (lines 947--950); eq:fbcontrol (lines 961--964); eq:Z-alg (lines 967--971); prop:sol2p1 (lines 1066--1071); eq:X-lin (lines 1120--1124); eq:Z-lin (lines 1120--1124); eq:tilde-Z-lin (lines 1120--1124); eq:approx-k-linear (lines 1146--1149). Every definition, display and label of those lines is retained unchanged. Omitted from this package and not redistributed: 1--676, 680--682, 686--787, 791--946, 951--960, 965--966, 972--1065, 1072--1119, 1125--1145, 1150--1150, 1178--1794. Nothing inside a retained excerpt is omitted or altered: every retained body is delivered exactly as the article's lines are written, so no omission receipt of any kind accompanies this package. Citations: the retained excerpts carry no citation command at all, so no bibliography entry of the article is transcribed and no reference is redistributed. Reference adaptation: none was needed and none was made; every reference inside the delivered unit resolves inside it, so no unresolved reference or citation is produced. Printed slips kept literal: no printed word, symbol, hypothesis or number of the counted statement or of its proof was altered. Layout and class: the article's own class is \texttt{ieeeconf} with packages including times, amsmath, amsfonts, amssymb, tabularx, ulem, dsfont, graphicx, subcaption, xcolor, hyperref, algorithm, algorithmic, float; the wrapper is the lane's pinned \texttt{amsart} editorial wrapper at 10pt on A4, which restates the theorem-like environments the retained text needs and adds the article's own macro definitions that the retained text needs (\texttt{\textbackslash Expect}, \texttt{\textbackslash R}, \texttt{\textbackslash trace}, \texttt{\textbackslash ud}), with extra wrapper packages (bm, bbm, dsfont, xspace, cancel, mathtools). The delivered numbering is the unit's own. Assets: no retained span contains a figure environment, a graphics-inclusion command, a PDF-page inclusion, a table or a TikZ picture; no image file is copied into this package, the package declares no asset dependency and no image file is redistributed with it. Licence: version arXiv:2504.00238v2 of this article is distributed under the Creative Commons Attribution 4.0 International licence, as recorded in the arXiv licence metadata for this exact version and shown on the abstract page https://arxiv.org/abs/2504.00238v2. A full rights notice is delivered at licenses/arxiv-2504.00238-CC-BY-4.0.txt. Characters: every retained excerpt is pure ASCII, so the delivered engine, XeLaTeX with its default Latin Modern text font, reports no missing character.
\begin{problem}\label{prob:steer-to-y}
    Given a target state $x_f \in  \R^n$ and a fixed terminal time $T>0$, find a feedback control law $\{U_t=k(t,X_t);t\in[0,T]\}$, for a function $k: [0,T] \times \R^n \rightarrow \R^m$, such that  $X_T = x_f$ almost surely.
\end{problem}

\begin{problem}\label{prob:approx-steer-to-y}
    Given a target state $x_f \in  \R^n$, a fixed terminal time $T>0$, and error tolerance $\delta>0$, find a feedback control law $\{U_t=k(t,X_t);t\in[0,T]\}$, for a function $k: [0,T] \times \R^n \rightarrow \R^m$, such that  $\mathbb E[\|X_T - x_f\|^2]\leq \delta$.
\end{problem} 

\begin{align}\label{eq:kappa-identity}
   &\kappa_{T,T-t}(x'|x) p(T-t,x) = \tilde \kappa_{t,0}(x|x') p(T,x')
\end{align}

\begin{align}\label{eq:k-star}
     k_i^\star(t,x) &\mathrel{:=} \frac{1}{ p(t,x)} \sum_{j=1}^n \partial_{x_j}(g_{j,i}(x)  p(t,x)),\\
 \mathfrak g_i(t,x) &\mathrel{:=} \sum_{j=1}^n \sum_{k=1}^m g_{j,k}(x) \partial_{x_j}g_{i,k}(x), \label{eq:frak-g}  
\end{align}

\begin{align}\label{eq:tilde-f}
    h(x) &:= -f(x) +  \epsilon^2 \mathfrak  g(x),\\
    k(t,x) &:= \epsilon^2 k^\star(T-t,x),\label{eq:fbcontrol}
\end{align}

\begin{align}\label{eq:Z-alg}
\ud Z_t &= (-f(Z_t) + \epsilon^2 \mathfrak g(Z_t))\ud t +  \epsilon g(Z_t) \ud W_t,~Z_0=x_f
\\ \ud \tilde Z_t &= 
    f(\tilde Z_t)\ud t+  g(\tilde Z_t) (k(t,\tilde Z_t)\ud t + \epsilon \ud \tilde W_t).\label{eq:tilde-Z-alg}
\end{align}

\begin{proposition}\label{prop:sol2p1}
% Consider problem~\ref{prob:steer-to-y} for the stochastic control system~\eqref{eq:sde}. 
Let $p(t,x)$ denote the probability density function  of $Z_t$ defined according to~\eqref{eq:Z-alg} and define the control law $k(t,x)$ according to~\eqref{eq:fbcontrol} and~\eqref{eq:k-star}. If the initial condition $X_0=x_0$ satisfies $p(T,x_0)>0$, then, problem~\ref{prob:steer-to-y} is solved with the feedback control law $k(t,x)$. 
    % The control input $U_t = k(t,X_t)$ can bring any initial state $\tilde x$, s.t. $P (\tilde X_T=\tilde x) >0$, to target state $y$ which means $U_t$ is the solution to Problem \ref{prob:p1}.
    % Assuming the conditions of Proposition \ref{prop:tr-sde} are valid for SDE \eqref{eq:alg-tr-sde} and probability density $\tilde p_t$, then Problem 1 is solved with the feedback control law \eqref{eq:fbcontrol}, if $X_0=x$ with $P(\tilde X_T = x)>0$.   
\end{proposition}

\begin{align}
    &\ud Z_t = -A Z_t \ud t + \epsilon B \ud W_t, \hspace{59pt} Z_0 = x_f,\label{eq:Z-lin}\\
    &\ud \tilde Z_t = A \tilde Z_t \ud t + B( k(t,\tilde Z_t) \ud t +\epsilon \ud \tilde W_t), \hspace{2pt} \tilde Z_T = x_f, \label{eq:tilde-Z-lin}\\
    &\ud X_t \!=\! A X_t \ud t \!+\! B(k(t,X_t)\ud t \!+ \!\epsilon \ud W_t), \hspace{10pt} X_0\! =\! x_0. \label{eq:X-lin}
\end{align}

\begin{align}\label{eq:approx-k-linear}
    \hat k(t,x) 
    &= -\epsilon^2 B^\top Q_{T-t}^{-1}(x-m_{T-t}).
\end{align}

\begin{proposition}\label{prop:terminal-error}
    % \at{Express the result (16) as a proposition and re-write the derivations as the proof. Express the fact that in the limit $\sigma \to 0$, the error goes to zero, as the second part of the proposition.}
    Under the feedback control law defined in \eqref{eq:approx-k-linear}, the expected squared error between the terminal state and the target state is given by
    \begin{equation}\label{eq:error}
        \Expect[\|X_T-x_f\|^2] = \sigma^4\|e^{TA}x_0-x_f\|^2_{M^{-2}} + \sigma^2(n -\trace(M^{-1})),
    \end{equation}
    where $M := \sigma^2 I +  \epsilon^2\int_0^T e^{As}BB^\top e^{A^\top s}\ud s$. Moreover, for any $\delta>0$, there exists a small enough $\sigma>0$ that solves Problem~\ref{prob:approx-steer-to-y}.
\end{proposition}

\begin{proof}
     Let $\kappa$ and $\tilde \kappa$ denote the probability transition kernels associated with SDEs~\eqref{eq:Z-lin} and~\eqref{eq:tilde-Z-lin}, respectively, similar to the proof of Proposition~\ref{prop:sol2p1}. In terms of the kernels, the probability distribution of $X_T$ is equal to $\tilde \kappa_{T,0}(\cdot|x_0)$, because the SDE~\eqref{eq:X-lin} and~\eqref{eq:tilde-Z-lin} have the same form and $X_0=x_0$.  Then, upon the application of the time-reversal relationship~\eqref{eq:kappa-identity}, and the fact that $p(t,x) = \mathcal N(x;m_t,Q_t)$,
\begin{align*}
    \tilde \kappa_{T,0}(x|x_0) &=  \frac{ \kappa_{T,0}(x_0|x)p(0,x)}{p(T,x_0)} \\
    & = \frac{\mathcal N(x_0;e^{-AT}x,\Sigma_T) \mathcal N(x;x_f,\sigma^2I)}{\mathcal N(x_0;m_T,Q_T)}\\&=\mathcal N(x;\mu,P)
\end{align*}
where 
\begin{align*}
    \mu &= x_f + \sigma^2 M^{-1}e^{TA}(x_0 - e^{-TA}x_f),\\
    P & =  \sigma^2(I- \sigma^2 M^{-1}).
\end{align*}
  Now, to compute the error   $\mathbb E[\|X_T-x_f\|^2]$, we use the fact that $X_T \sim \mathcal N(\cdot;\mu, P)$ and, hence: 
\begin{align}
    \mathbb E[\|X_T-x_f\|^2] &= \|\mu-x_f\|^2 + \trace(P) \nonumber
    % \\
    % &= \sigma^4\|e^{TA}x_0\!-\!x_f\|^2_{M^{-2}} \!+\! \sigma^2(n\! -\!\trace(M^{-1})), \nonumber
\end{align}
which yields~\eqref{eq:error}. Moreover, in the limit as $\sigma \to 0$, $M$~converges to the controllability grammarian matrix which is positive definite  under the controllability assumption.  Therefore, taking the limit of~\eqref{eq:error} as $\sigma \to 0$, concludes that the error converges to zero, implying that for any $\delta>0$, there exists a $\sigma>0$ such that the error is smaller than $\delta$. 
\end{proof}

\end{document}
